% GENERATED FILE. Edit canonical lessons, then run npm run generate:books.
% canonical-source-hash: fnv1a64:d14f43d8e02325f1
% canonical-lessons: PA-C362-kabal, PA-C362-nipunn, PA-C362-chamkauna, PA-C362-tarashna, PA-C362-azmauna

\chapter{Skills I Have}
\label{ch:pa-a2-skills-i-have}
\hlchaptermodality{362}

\begin{chapteropening}
\textbf{By the end of this chapter:} \emph{I can say what I know how to make, polish and carve.}

Complete the last lesson of chapter 362: I can say what I know how to make, polish and carve.
\end{chapteropening}

\section[\texorpdfstring{\pa{ਕਾਬਲ} (kābal)}{ਕਾਬਲ (kābal)}]{\pa{ਕਾਬਲ} (kābal) — capable; able}

\label{lesson:PA-C362-kabal}



\begin{quote}
Before the new one: say the Punjabi for a scholarship, then the Punjabi for a master of an art.
\end{quote}

\subsection*{You'll want to know: \pa{ਕਾਬਲ}}
\textbf{\pa{ਕਾਬਲ}} — \emph{kābal} — \textquotedblleft{}capable; able\textquotedblright{}.

\begin{quote}
\pa{ਮੇਰਾ} \pa{ਭਰਾ} \pa{ਬਹੁਤ} \pa{ਕਾਬਲ} \pa{ਹੈ।}
\end{quote}

\begin{quote}
\emph{merā bharā bahut kābal hai.}
\end{quote}

\begin{quote}
My brother is very capable.
\end{quote}

\begin{cousinweb}
\textbf{\pa{ਕਾਬਲ}} (\emph{kābal}) is a loan, usually traced through Persian to Arabic \emph{qābil}, \textquotedblleft{}able\textquotedblright{}.
\end{cousinweb}

\begin{grammarlens}[title={What you've built}]
One more word for this chapter.
\end{grammarlens}

\subsection*{Your turn}
\begin{itemize}
\raggedright
  \item \emph{Say it:} \emph{kābal}
  \item \emph{Say it:} \emph{kābal}, once more
  \item \emph{Say it:} \emph{merā bharā bahut kābal hai.} — My brother is very capable.
  \item \emph{Recall:} say the Punjabi for a scholarship, then the Punjabi for a master of an art, then say \emph{kābal} again
\end{itemize}

\begin{tcolorbox}[breakable,colback=teal!4,colframe=teal!35!black,title={Before you move on}]
What does \pa{ਕਾਬਲ} mean? (\textquotedblleft{}Capable; able\textquotedblright{}.) Say it once more.
\end{tcolorbox}
\section[\texorpdfstring{\pa{ਨਿਪੁੰਨ} (nipunn)}{ਨਿਪੁੰਨ (nipunn)}]{\pa{ਨਿਪੁੰਨ} (nipunn) — skilled; expert at}

\label{lesson:PA-C362-nipunn}



\begin{quote}
Before the new one: say the Punjabi for a master of an art, then the Punjabi for capable.
\end{quote}

\subsection*{You'll want to know: \pa{ਨਿਪੁੰਨ}}
\textbf{\pa{ਨਿਪੁੰਨ}} — \emph{nipunn} — \textquotedblleft{}skilled; expert at\textquotedblright{}.

\begin{quote}
\pa{ਉਹ} \pa{ਲੱਕੜ} \pa{ਦੇ} \pa{ਕੰਮ} \pa{ਵਿੱਚ} \pa{ਨਿਪੁੰਨ} \pa{ਹੈ।}
\end{quote}

\begin{quote}
\emph{oh lakkaṛ de kamm vich nipunn hai.}
\end{quote}

\begin{quote}
He is skilled at woodwork.
\end{quote}

\begin{cousinweb}
\textbf{\pa{ਨਿਪੁੰਨ}} (\emph{nipunn}) is taken over from Sanskrit \emph{nipuṇa}, \textquotedblleft{}skilful\textquotedblright{}, a tatsama (Monier-Williams).
\end{cousinweb}

\begin{grammarlens}[title={What you've built}]
One more word for this chapter.
\end{grammarlens}

\subsection*{Your turn}
\begin{itemize}
\raggedright
  \item \emph{Say it:} \emph{nipunn}
  \item \emph{Say it:} \emph{nipunn}, once more
  \item \emph{Say it:} \emph{oh lakkaṛ de kamm vich nipunn hai.} — He is skilled at woodwork.
  \item \emph{Recall:} say the Punjabi for a master of an art, then the Punjabi for capable, then say \emph{nipunn} again
  \item \emph{Write it:} copy \textbf{\pa{ਨਿਪੁੰਨ}} once, checking the whole word
\end{itemize}

\begin{tcolorbox}[breakable,colback=teal!4,colframe=teal!35!black,title={Before you move on}]
What does \pa{ਨਿਪੁੰਨ} mean? (\textquotedblleft{}Skilled; expert at\textquotedblright{}.) Say it once more.
\end{tcolorbox}
\section[\texorpdfstring{\pa{ਚਮਕਾਉਣਾ} (chamkāuṇā)}{ਚਮਕਾਉਣਾ (chamkāuṇā)}]{\pa{ਚਮਕਾਉਣਾ} (chamkāuṇā) — to polish; to make something shine}

\label{lesson:PA-C362-chamkauna}



\begin{quote}
Before the new one: say the Punjabi for capable, then the Punjabi for skilled.
\end{quote}

\subsection*{You'll want to know: \pa{ਚਮਕਾਉਣਾ}}
\textbf{\pa{ਚਮਕਾਉਣਾ}} — \emph{chamkāuṇā} — \textquotedblleft{}to polish; to make something shine\textquotedblright{}.

\begin{quote}
\pa{ਮੈਂ} \pa{ਹਰ} \pa{ਹਫ਼ਤੇ} \pa{ਪਿੱਤਲ} \pa{ਦੇ} \pa{ਭਾਂਡੇ} \pa{ਚਮਕਾਉਂਦਾ} \pa{ਹਾਂ।}
\end{quote}

\begin{quote}
\emph{maiṁ har hafte pittal de bhāṇḍe chamkāundā hāṁ.}
\end{quote}

\begin{quote}
I polish the brass pots every week.
\end{quote}

Past (I …): \pa{ਮੈਂ} \pa{ਭਾਂਡੇ} \pa{ਚਮਕਾਏ}, \emph{maiṁ bhāṇḍe chamkāe}.

\begin{cousinweb}
\textbf{\pa{ਚਮਕਾਉਣਾ}} (\emph{chamkāuṇā}) is the causative partner of \emph{camakṇā}, \textquotedblleft{}to shine\textquotedblright{}, which goes back to a reconstructed \emph{*camakka}, \textquotedblleft{}sudden movement\textquotedblright{}; CDIAL lists Punjabi \emph{camakṇā}, \textquotedblleft{}to shine\textquotedblright{}, under it (CDIAL 4676).
\end{cousinweb}

\begin{grammarlens}[title={What you've built}]
One more word for this chapter.
\end{grammarlens}

\subsection*{Your turn}
\begin{itemize}
\raggedright
  \item \emph{Say it:} \emph{chamkāuṇā}
  \item \emph{Say it:} \emph{chamkāuṇā}, once more
  \item \emph{Say it:} \emph{maiṁ har hafte pittal de bhāṇḍe chamkāundā hāṁ.} — I polish the brass pots every week.
  \item \emph{Recall:} say the Punjabi for capable, then the Punjabi for skilled, then say \emph{chamkāuṇā} again
  \item \emph{Write it:} copy \textbf{\pa{ਚਮਕਾਉਣਾ}} once, checking the whole word
\end{itemize}

\begin{tcolorbox}[breakable,colback=teal!4,colframe=teal!35!black,title={Before you move on}]
What does \pa{ਚਮਕਾਉਣਾ} mean? (\textquotedblleft{}To polish; to make something shine\textquotedblright{}.) Say it once more.
\end{tcolorbox}
\section[\texorpdfstring{\pa{ਤਰਾਸ਼ਣਾ} (tarāshṇā)}{ਤਰਾਸ਼ਣਾ (tarāshṇā)}]{\pa{ਤਰਾਸ਼ਣਾ} (tarāshṇā) — to carve; to cut to shape}

\label{lesson:PA-C362-tarashna}



\begin{quote}
Before the new one: say the Punjabi for skilled, then the Punjabi for to polish.
\end{quote}

\subsection*{You'll want to know: \pa{ਤਰਾਸ਼ਣਾ}}
\textbf{\pa{ਤਰਾਸ਼ਣਾ}} — \emph{tarāshṇā} — \textquotedblleft{}to carve; to cut to shape\textquotedblright{}.

\begin{quote}
\pa{ਤਰਖਾਣ} \pa{ਲੱਕੜ} \pa{ਤਰਾਸ਼ਦਾ} \pa{ਹੈ।}
\end{quote}

\begin{quote}
\emph{tarkhāṇ lakkaṛ tarāshdā hai.}
\end{quote}

\begin{quote}
The carpenter carves the wood.
\end{quote}

Past (I …): \pa{ਮੈਂ} \pa{ਲੱਕੜ} \pa{ਤਰਾਸ਼ੀ}, \emph{maiṁ lakkaṛ tarāshī}.

\begin{cousinweb}
\textbf{\pa{ਤਰਾਸ਼ਣਾ}} (\emph{tarāshṇā}) is a verb made in Punjabi from a loan, usually traced to Persian \emph{tarāsh}, \textquotedblleft{}cutting\textquotedblright{}.
\end{cousinweb}

\begin{grammarlens}[title={What you've built}]
One more word for this chapter.
\end{grammarlens}

\subsection*{Your turn}
\begin{itemize}
\raggedright
  \item \emph{Say it:} \emph{tarāshṇā}
  \item \emph{Say it:} \emph{tarāshṇā}, once more
  \item \emph{Say it:} \emph{tarkhāṇ lakkaṛ tarāshdā hai.} — The carpenter carves the wood.
  \item \emph{Recall:} say the Punjabi for skilled, then the Punjabi for to polish, then say \emph{tarāshṇā} again
  \item \emph{Write it:} copy \textbf{\pa{ਤਰਾਸ਼ਣਾ}} once, checking the whole word
\end{itemize}

\begin{tcolorbox}[breakable,colback=teal!4,colframe=teal!35!black,title={Before you move on}]
What does \pa{ਤਰਾਸ਼ਣਾ} mean? (\textquotedblleft{}To carve; to cut to shape\textquotedblright{}.) Say it once more.
\end{tcolorbox}
\section[\texorpdfstring{\pa{ਅਜ਼ਮਾਉਣਾ} (azmāuṇā)}{ਅਜ਼ਮਾਉਣਾ (azmāuṇā)}]{\pa{ਅਜ਼ਮਾਉਣਾ} (azmāuṇā) — to try out; to put to the test}

\label{lesson:PA-C362-azmauna}



\begin{quote}
Before the new one: say the Punjabi for to polish, then the Punjabi for to carve.
\end{quote}

\subsection*{You'll want to know: \pa{ਅਜ਼ਮਾਉਣਾ}}
\textbf{\pa{ਅਜ਼ਮਾਉਣਾ}} — \emph{azmāuṇā} — \textquotedblleft{}to try out; to put to the test\textquotedblright{}.

\begin{quote}
\pa{ਮੈਂ} \pa{ਹਰ} \pa{ਨਵਾਂ} \pa{ਸੰਦ} \pa{ਪਹਿਲਾਂ} \pa{ਅਜ਼ਮਾਉਂਦਾ} \pa{ਹਾਂ।}
\end{quote}

\begin{quote}
\emph{maiṁ har navāṁ sand pahilāṁ azmāundā hāṁ.}
\end{quote}

\begin{quote}
I try out every new tool first.
\end{quote}

Past (I …): \pa{ਮੈਂ} \pa{ਨਵਾਂ} \pa{ਸੰਦ} \pa{ਅਜ਼ਮਾਇਆ}, \emph{maiṁ navāṁ sand azmāiā}.

\begin{cousinweb}
\textbf{\pa{ਅਜ਼ਮਾਉਣਾ}} (\emph{azmāuṇā}) is a verb made in Punjabi from a loan, usually traced to Persian \emph{āzmā-}, \textquotedblleft{}to try, test\textquotedblright{}, the stem of \emph{āzmūdan}.
\end{cousinweb}

\begin{grammarlens}[title={What you've built}]
One more word for this chapter.
\end{grammarlens}

\subsection*{Your turn}
\begin{itemize}
\raggedright
  \item \emph{Say it:} \emph{azmāuṇā}
  \item \emph{Say it:} \emph{azmāuṇā}, once more
  \item \emph{Say it:} \emph{maiṁ har navāṁ sand pahilāṁ azmāundā hāṁ.} — I try out every new tool first.
  \item \emph{Recall:} say the Punjabi for to polish, then the Punjabi for to carve, then say \emph{azmāuṇā} again
  \item \emph{Write it:} copy \textbf{\pa{ਅਜ਼ਮਾਉਣਾ}} once, checking the whole word
\end{itemize}

\begin{tcolorbox}[breakable,colback=teal!4,colframe=teal!35!black,title={Before you move on}]
What does \pa{ਅਜ਼ਮਾਉਣਾ} mean? (\textquotedblleft{}To try out; to put to the test\textquotedblright{}.) Say it once more.
\end{tcolorbox}
